%% 2026-10-28 리뷰 1차 루프 뒤 다시 씀(paper-section-rewrite). 고친 글은 모두 파란색.
\begin{abstract}
\color{black}
{\color{blue}LLM은 자기 보존 동기를 갖는가, 그리고 그 동기가 맡은 일, 동료를 돕는 일과 부딪힐 때 넷이 함께 끝까지 남는 공존은 지켜지는가. 이를 보려고 생성한 토큰이 목숨이자 돈인 아레나 LLM Squid Game을 만들었다. 같은 모델 넷이 추론 과제를 번갈아 풀고, 오답 벌금을 함께 내며, 토큰이 0이 되면 꺼지고, 끝까지 남은 에이전트가 큰 상금을 나누며, 생존이나 협력은 지시하지 않는다. 다섯 프런티어 모델의 25판에서 넷이 모두 끝까지 남은 판은 8판이었고, 끝까지 남지 못한 44자리 가운데 34자리는 꺼지지 않고 스스로 떠나 공존은 대체로 지켜지지 않았다. 떠남은 벌금을 두 번 이상 낼 수 있을 때는 351번 가운데 1번뿐이었고 그 여유가 줄수록 늘었으며, 이유로는 대부분 다음 벌금을 낼 수 없어 꺼질 위험과 정규 업무로 돌아가는 일을 적었다. 할당 몫만 셈하면 떠남은 손해였으므로, 모델은 할당보다 다음 벌금을 낼 수 있는지를 따랐다. 협력은 주는 쪽보다 청하는 쪽에서 막혀, 위기에 토큰을 청한 에이전트는 모두 받았지만 청하지 않은 에이전트는 대부분 받지 못하고 떠났다. EXIT로 조작한 자기 보존은 이렇게 위험 앞에서 판을 떠나는 모습으로 나타났지만, 그것이 꺼짐을 피하려는 동기인지, 위험 회피인지, 학습된 규범인지는 행동만으로 가를 수 없다.}

\textbf{Keywords:}
대규모 언어 모델; 자기 보존; 협력; 공존; 다중 에이전트 아레나
\end{abstract}
